\documentclass[12pt, letterpaper]{article}
\usepackage[english]{babel}
\usepackage[T1]{fontenc}
\usepackage[margin=1.15in]{geometry}
\usepackage{xcolor}
\usepackage{url}
\usepackage[utf8]{inputenc}
\usepackage{tabularx}
\usepackage{booktabs}
\frenchspacing
\usepackage{multicol}
\usepackage{eso-pic}
\usepackage[longnamesfirst]{natbib}
\bibpunct{(}{)}{;}{a}{}{,}
\usepackage{caption}
\usepackage{subcaption}
\usepackage{setspace}
\usepackage{paralist}
\usepackage{quoting}
\usepackage{comment}
\usepackage{enumitem}
\usepackage{graphicx}
\usepackage{float}
\usepackage{bookmark}
\renewcommand{\thesection}{\arabic{section}.}
\renewcommand{\thesubsection}{\thesection\arabic{subsection}}
\renewcommand{\thesubsubsection}{\thesubsection.\arabic{subsubsection}}
\usepackage{hyperref}
\hypersetup{
    colorlinks=true,
        linkcolor=blue,
        filecolor=magenta,      
        urlcolor=cyan,
        citecolor=purple,
}
\usepackage[scaled]{roboto} % Load the helvet package
\renewcommand*\familydefault{\sfdefault} % Set the default font to be sans-serif
\newenvironment{boldnumlist}
{\begin{enumerate}[label=\textbf{\arabic*.}]}
{\end{enumerate}}

\begin{document}
\title{\includegraphics[width=8cm]{Images/stacked.png} \\ \textbf{Introduction to Public Administration}}

\title{MPA Capstone Seminar: \\ Public Administration Theory}
\author{POSC 521 — Spring 2025}
\date{Tuesdays at 7:00 in GH 208}
    \maketitle
        \subsection*{Session Schedule:}
                    
                    \begin{multicols}{2}
                    \subsubsection*{In-Person Sessions}
                    \begin{itemize}[leftmargin=*]
                        \item January 28
                        \item February 4, 11, 18, 25
                        \item March 4, 11, 18, 25
                        \item May 6
                    \end{itemize}
                    
                    \columnbreak
                    
                    \subsubsection*{Asynchronous Online Sessions}
                    \begin{itemize}[leftmargin=*]
                        \item January 21
                        \item April 8, 15, 22, 29
                        \item May 13
                    \end{itemize}
                    \end{multicols}


                    \subsection*{Professor: David P. Adams, Ph.D.}

                    \subsubsection*{Contact Information:}
                    
                    \begin{itemize}
                            \item Office: 516 Gordon Hall
                            \item Phone/SMS: (657) 278-4770
                            \item Zoom Meeting ID: 334 750 2369 or \href{https://fullerton.zoom.us/j/3347502369}{\texttt{fullerton.zoom.us/j/3347502369}} 
                            \item website: \href{https://dadams.io}{\texttt{dadams.io}}
                            \item email: \href{dpadams@fullerton.edu}{\texttt{dpadams@fullerton.edu}}
                            \item Office hours: Tuesdays \& Thursdays from 9:30 to 11:00, Thursdays from 5:30 to 6:30, and by \href{https://dadams.io/appointments}{appointment}.
                            \item Schedule meetings throughout the week: \href{https://dadams.io/appointments}{\texttt{dadams.io/appointments}}
                    \end{itemize}
                    
                    \section{Catalog Description}
                    Concepts, models and ideologies of public administration within the larger political system. Course restricted to students in their final six units of graduate work.
                    
                    \section{Course Description}
                    The capstone seminar in the Master of Public Administration program at Cal State Fullerton examines concepts, models, and ideologies of public administration within the larger political system.

                    \section{Course Prerequisites}
                    This course is restricted to students in their final six units of graduate work in the MPA program. Students must have completed all other required courses in the MPA program before enrolling in this course.
                    
                    \section{Course Objectives}
                    This course is designed to accomplish five interrelated objectives:
                    
                    \begin{enumerate}
                            \item \textbf{Theory Examination}: We will delve into the most important theories and literature in public administration, fostering a deep understanding of the field.
                            \item \textbf{Literature Review}: You will complete a literature review in your concentration area, allowing you to specialize and delve deeper into a specific aspect of public administration. This preparation will be crucial for the general concentration portion of the comprehensive exams.
                            \item \textbf{Writing Skills}: This course will enhance your writing skills, focusing on clear, concise, and effective communication. This preparation will be crucial for the general theory portion of the comprehensive exams.
                            \item \textbf{Critical Thinking}: You will develop your critical thinking skills, learning to analyze and evaluate complex arguments and theories. This preparation will be crucial for the general theory portion of the comprehensive exams.
                            \item \textbf{Professional Development}: This course will help you develop the skills and knowledge necessary for a successful career in public administration. You will learn about the latest trends and issues in the field and how to navigate the challenges of public service.
                    
                    \end{enumerate}
    
\section{University-wide Student Learning Outcomes}
    As a capstone course in the MPA program, this course is designed to help students achieve the following university-wide student learning outcomes:
    \begin{enumerate}
        \item Knowledge, skills, and professional dispositions including higher order competence in disciplinary perspectives and interdisciplinary points of view;
        \item The ability to access, analyze, synthesize, and evaluate complex information from multiple sources and in new situations and settings;
        \item Advanced communication skills;
        \item The ability to work independently and in collaboration with others as artists, practitioners, researchers, and/or scholars;
        \item The ability to determine and apply appropriate methods and technologies to address problems that affect their communities; 
        \item A commitment to social justice and ethical leadership within diverse communities and an interdependent global community.
    \end{enumerate}

    \section{Course Materials}
    \subsection*{Required Texts}
    \begin{itemize}
        \item \textbf{Denhardt and Denhardt}. \textit{The New Public Service: Serving, Not Steering}. 4th ed. Routledge, 2015.
        \item \textbf{Hanna-Attisha, Mona}. \textit{What the Eyes Don't See: A Story of Crisis, Resistance, and Hope in an American City}. One World, 2018.
        \item \textbf{Lipsky, Michael}. \textit{Street-Level Bureaucracy: Dilemmas of the Individual in Public Services}. Russell Sage Foundation, 2010.
        \item \textbf{Stivers, Camilla}. \textit{Bureau Men, Settlement Women: Constructing Public Administration in the Progressive Era}. University of Kansas Press, 2000.
    \end{itemize}

\subsection*{Additional Readings}

\noindent Additional readings are indicated on the course schedule and in the references section of this document. These readings and a \texttt{.bib} file are located in a shared DropBox folder. Please contact the professor if you need access to the folder.

\section{Technical Requirements}

\subsection*{Pollak Library Resources}

The Pollak Library provides a wide range of resources and services to support your research and learning. These resources include books, journals, databases, and research guides. You can access the library's resources online through the \href{http://www.library.fullerton.edu/}{Pollak Library website}. The library also offers research assistance through the \href{http://www.library.fullerton.edu/research/}{Research Assistance Program}. You can also access the \href{http://www.library.fullerton.edu/about/guidelines/online-instruction-guidelines.php}{library's online instruction guidelines} for help with online learning.

\subsection*{Canvas}

This course will use \href{https://csufullerton.instructure.com/}{Canvas} as a learning management system. You will use \emph{Canvas} to access course materials, submit assignments, participate in discussions, and communicate with the professor and your classmate. You are responsible for checking \emph{Canvas} regularly for announcements, assignments, and other course materials. You are also responsible for ensuring that your \emph{Canvas} notifications are set to receive messages from the course. 

\subsection*{Zoom}
This course may include synchronous online sessions using \href{https://fullerton.zoom.us/}{Zoom}. You are responsible for ensuring that you have the necessary equipment and internet connection to participate in these sessions. 

\subsection*{Minimum Technical Requirements}

To participate in this course, you will need the following minimum technical requirements:
\begin{itemize}
    \item A computer or tablet with a reliable internet connection
    \item A webcam and microphone
    \item A modern web browser (Chrome, Firefox, Safari, or Edge)
    \item Microsoft Word or a compatible word processing program
    \item Adobe Acrobat Reader or a compatible PDF reader
\end{itemize}


\noindent Long- and short-term computer and internet access loans are available through the \href{http://www.fullerton.edu/it/students/sgc/index.php}{Student Genius Center}.

\section{Technical Problems}

\subsection*{University IT Help Desk}

Contact the instructor immediately to document the problem if you encounter any technical difficulties. Then contact the \href{http://www.fullerton.edu/it/students/helpdesk/index.php}{Student IT Help Desk} for assistance. You can also call the Student IT Help Desk at (657) 278-8888, \href{mailto:StudentITHelpDesk@fullerton.edu}{email}, visit them at the Pollak Library North \href{http://www.fullerton.edu/it/students/sgc/index.php}{Student Genius Center}, or log on to the \href{http://my.fullerton.edu/}{my.fullerton.edu} portal and click ``Online IT Help'' followed by ``Live Chat''.

\subsection*{Canvas Support}

If you encounter any technical difficulties with Canvas, call the Canvas Support Hotline at 855-302-7528, visit the \href{https://community.canvaslms.com/docs/DOC-10720-67952720329}{Canvas Community}, or click the ``Help'' button in the lower left corner of Canvas and select ``Report a Problem''. The \href{https://cases.canvaslms.com/liveagentchat?chattype=student&sfid=001A000000YzcwQIAR}{Student Support Live Chat} is available 24 hours a day, 7 days a week.


\section{University Student Policies}

In accordance with UPS 300.00, students must be familiar with certain policies applicable to all courses. Please review these policies as needed and visit this Cal State Fullerton website \texttt{\href{https://fdc.fullerton.edu/teaching/student-info-syllabi.html}{https://fdc.fullerton.edu/teaching/student-info-syllabi.html}} for links to the following information:

\begin{enumerate}
    \item   University learning goals and program learning outcomes.
    \item	Learning objectives for each General Education (GE) category.
    \item	Guidelines for appropriate online behavior (netiquette).
    \item	Students' rights to accommodations for documented special needs.
    \item   Campus student support measures, including Counseling \& Psychological Services, Title IV and Gender Equity, Diversity Initiatives and Resource Centers, and Basic Needs Services.
    \item	Academic integrity (refer to UPS 300.021).
    \item	Actions to take during an emergency.
    \item	Library services information.
    \item	Student Information Technology Services, including details on technical competencies and resources required for all students.
    \item	Software privacy and accessibility statements.
\end{enumerate}

\section{Course Student Policies}

\subsection*{Course Communication}
All course announcements and communications will be sent via \emph{Canvas} and university email. Students are responsible for regularly checking their \emph{Canvas} notifications and email. Students are also responsible for ensuring that their \emph{Canvas} notifications are set to receive messages from the course. Students are expected to check \emph{Canvas} and their email at least once daily.

\subsubsection*{Response Time}I will strive to respond to all student emails and \emph{Canvas} messages within 24 hours, except on weekends and holidays. If you do not receive a response within 24 hours, please send a follow-up message. If you do not receive a response within 48 hours, please send another follow-up message and contact me via phone or SMS text at (657) 278-4770.

\subsection*{Due Dates}
All assignments are due on the date specified in the course schedule. Late assignments will only be accepted if prior arrangements have been made with the professor. Students must submit all assignments on time and in the correct format. Failure to submit an assignment on time may result in a grade penalty.

\subsection*{Alternative Procedures for Submitting Work}
Students are expected to submit all assignments via \emph{Canvas}. If you cannot submit an assignment via \emph{Canvas}, please contact the professor to discuss alternative submission procedures.

\subsection*{Extra Credit}
Extra credit opportunities will not be offered in this course. All students will be graded based on the same criteria and standards.

\subsection*{Attendance}
Students are expected to attend all in-person sessions. If you are unable to attend a session, please notify the professor in advance. If you miss a session, you are responsible for obtaining the information and materials covered in the session.

\subsection*{Retention of Student Work}
Students are responsible for retaining copies of all assignments submitted in this course. Students should keep copies of all assignments until the end of the semester and verify that their assignments have been graded and returned before discarding them.

\subsection*{Academic Integrity}
Students are expected to adhere to the highest standards of academic integrity. Any student found to have engaged in academic dishonesty will be subject to the sanctions described in the \href{https://www.fullerton.edu/senate/publications_policies_resolutions/ups/UPS%20300/UPS%20300.021.pdf}{Academic Dishonesty Policy} (UPS 300.021). Academic dishonesty includes, but is not limited to, cheating, plagiarism, fabrication, facilitating academic dishonesty, and submitting previously graded work without prior authorization. Students are expected to be familiar with the university's policy on academic dishonesty and to adhere to this policy in all aspects of this course. Any student who has questions about the policy should ask the professor for clarification.

\subsection*{Plagiarism}
Plagiarism is a serious violation of academic integrity and will not be tolerated in this course. Plagiarism includes, but is not limited to, copying and pasting text from sources without proper citation, paraphrasing text from sources without proper citation, and submitting work that is not your own. Students are expected to properly cite all sources used in their work and to submit original work. Failure to do so may result in a failing grade for the assignment and further disciplinary action.

\subsection*{Written Work}
All written work must be submitted in a professional format, including proper grammar, spelling, and punctuation. Written work must also be properly cited using the appropriate citation style. Students are expected to follow the guidelines for written work provided by the professor and to seek clarification if they have questions about the requirements.

\subsection*{Artificial Intelligence Policy}
\subsubsection*{Definitions of Generative AI}

\noindent For the purposes of this course, generative AI refers to artificial intelligence systems capable of producing human-like text, images, or other content. This includes, but is not limited to:

\begin{itemize}
    \item Large language models (LLM) (e.g., GPT-3, GPT-4)
    \item Text-to-image generators (e.g., DALL-E, Midjourney)
    \item AI-powered writing assistants (e.g., ChatGPT, Claude)
    \item Automated content generators
\end{itemize}

\subsubsection*{AI Use Policy}

In this course, the use of AI is permitted and encouraged as a learning tool, with the following guidelines:

\begin{itemize}
    \item \textbf{Allowed}: Using AI for feedback on annotated bibliographies and syntheses, as outlined in the assignment.
    \item \textbf{Allowed}: Using AI for the reflection activity.
    \item \textbf{Not Allowed}: Submitting AI-generated content as your own work without substantial human input and critical analysis.
    \item \textbf{Not Mandatory}: The use of AI is optional. Students are not required to use AI tools if they prefer not to.
\end{itemize}

\subsubsection*{Rationale for AI Policy}

This AI policy is designed to align with our course objectives:

\begin{enumerate}
    \item It supports the examination of public administration theories by encouraging critical engagement with AI feedback.
    \item It enhances literature review skills by allowing AI to provide additional perspectives on annotations and syntheses.
    \item It improves writing skills by providing an additional source of feedback and encouraging revision.
    \item It fosters critical thinking by requiring students to evaluate and incorporate AI feedback thoughtfully.
    \item It contributes to professional development by familiarizing students with AI tools they may encounter in their public administration careers.
    \item It promotes awareness of ethical considerations and sustainability impacts associated with AI use.
\end{enumerate}

\subsubsection*{Ethics and Sustainability Considerations}

\noindent Students are expected to engage with AI tools ethically and responsibly. This includes:

\begin{itemize}
    \item \textbf{Ethical Use}: Avoid using AI to plagiarize or produce content without proper attribution. Cite AI tools as sources when they contribute to your work.
    \item \textbf{Bias Awareness}: Understand that AI tools can reflect and perpetuate biases present in their training data. Critically evaluate the outputs for fairness and inclusivity.
    \item \textbf{Sustainability Impact}: Be mindful of the environmental impact of AI, including the carbon emissions associated with large-scale model training and inference. Use AI judiciously and consider alternative methods when feasible.
\end{itemize}

\subsubsection*{Guidance for Students on AI Use}

\begin{itemize}
    \item When using AI for feedback, always cite it as a source. For example: ``ChatGPT (Version X.X) [AI model]. OpenAI. Feedback received on [Date].''
    \item Be critical of AI-generated content. AI can make mistakes or provide biased information.
    \item Use AI as a tool to enhance your learning, not as a substitute for your own critical thinking and analysis.
    \item We will discuss ethical AI use, sustainability impacts, and its implications for public administration during our course.
\end{itemize}

\subsubsection*{Repercussions for Policy Breaches}

\begin{itemize}
    \item Suspected misuse of AI (e.g., submitting AI-generated content as your own without proper attribution) will be addressed on a case-by-case basis.
    \item Detection methods may include inconsistencies in writing style, unusual phrasing, or content that doesn't align with the student's demonstrated knowledge and skills.
    \item Consequences for policy breaches may range from a required revision of the assignment to more severe academic integrity penalties, such as dismissal from the MPA program, depending on the nature and extent of the misuse.
\end{itemize}

\noindent Remember, the goal of this policy is to enhance your learning experience and prepare you for a future where AI will likely play a significant role in public administration. Use these tools responsibly and ethically to support your academic and professional growth.

\subsubsection*{Goals and Assessment of AI Use}

\noindent The goals for AI use in this course are:

\begin{enumerate}
    \item To teach students to use AI as a complementary tool for enhancing critical thinking and writing skills.
    \item To prepare students for professional environments where AI tools are increasingly common in public administration.
    \item To encourage critical evaluation of AI-generated feedback and content.
    \item To foster ethical and responsible use of AI in academic and professional settings.
    \item To promote sustainability awareness by considering the environmental costs of AI use.
\end{enumerate}

\noindent Assessment of these goals and overall learning objectives will be conducted through:

\begin{itemize}
    \item Comparison of initial drafts with AI-assisted revisions
    \item Analysis of student reflections on the learning process
    \item Evaluation of class participation and discussions
    \item Periodic comparative assignments (with and without AI assistance)
    \item Review of AI usage logs
    \item Assessment of peer review quality
    \item Comprehensive final project or exam
    \item Professional development reflection on future AI use in public administration
\end{itemize}

\noindent These assessment methods will help ensure that AI use is enhancing, rather than replacing, critical thinking and analytical skills essential to public administration.

\section{Course Requirements}
\subsection*{Weekly Readings Assignment}

The capstone course is designed to integrate key theories and practices in public administration through weekly assignments culminating in a comprehensive capstone synthesis paper. Students will engage with foundational readings, participate in in-person discussions, and produce critical analyses that connect theory to practice.

\subsubsection*{Weekly Assignment Workflow}

Each week includes the following components:

\begin{enumerate}
    \item \textbf{Annotated Bibliography (Due Tuesday Before Class)}
    \begin{itemize}
        \item \textbf{Purpose}: Prepare for in-class discussion by summarizing key ideas and evaluating the relevance of assigned readings.
        \item \textbf{Instructions}:
        \begin{itemize}
            \item For each assigned reading:
            \begin{itemize}
                \item Provide a citation in APA or Chicago author-date style.
                \item Write a 150-word annotation that includes:
                \begin{itemize}
                    \item A brief summary of the central argument and key points.
                    \item Relevance to the week’s topic and implications for public administration.
                \end{itemize}
            \end{itemize}
            \item \textbf{Optional AI Integration}: Students may use Open AI's \href{https://chat.openai.com/}{ChatGPT}, Anthropic's \href{https://www.anthropic.com/}{Claude}, Google's \href{https://www.google.com/search/about/}{Gemini}, or the university's \href{https://titangpt.fullerton.edu/auth/jwt/login}{TitanGPT} to review their annotations:
            \begin{quote}
                \texttt{"I've completed an annotated bibliography for an article. The citation is [insert citation]. The summary I wrote is [insert summary]. I've noted the relevance as [insert relevance]. Could you provide feedback on my summary and relevance note?"}

            \end{quote}
        \end{itemize}
    \end{itemize}

    \item \textbf{In-Person Class Discussion (Tuesday)}
    \begin{itemize}
        \item \textbf{Purpose}: Deepen understanding of the week’s readings through collaborative discussion.
        \item \textbf{Activity}: Analyze key themes, patterns, and contradictions in the readings.
    \end{itemize}

    \item \textbf{Synthesis Paper (Due Thursday Night)}
    \begin{itemize}
        \item \textbf{Purpose}: Develop a critical analysis integrating insights from the readings and in-class discussion.
        \item \textbf{Instructions}:
        \begin{itemize}
            \item Write a 3- to 5-page synthesis paper that:
            \begin{itemize}
                \item Identifies patterns, connections, and contradictions across the readings.
                \item Highlights implications for public administration theory and practice.
            \end{itemize}
            \item \textbf{AI Integration}:
            \begin{quote}
                \texttt{"I've synthesized information from [number] articles on [topic]. Here's my draft of the synthesis [insert draft]. Could you provide feedback and suggest any connections or contrasts I might have missed?"}

            \end{quote}
        \end{itemize}
    \end{itemize}

    \item \textbf{Peer Review (Due Friday Night)}
    \begin{itemize}
        \item \textbf{Purpose}: Strengthen synthesis papers through collaborative feedback.
        \item \textbf{Instructions}:
        \begin{itemize}
            \item Review two classmates’ papers using a rubric provided on Canvas.
            \item Provide constructive feedback focusing on clarity, depth of analysis, and connections between readings.
        \end{itemize}
    \end{itemize}

    \item \textbf{Final Synthesis Paper \& Reflection (Due Saturday Night)}
    \begin{itemize}
        \item \textbf{Purpose}: Incorporate feedback to refine the synthesis paper.
        \item \textbf{Instructions}:
        \begin{itemize}
            \item Revise the synthesis paper based on peer feedback.
            \item Submit the final synthesis to Canvas.
        \end{itemize}
    \end{itemize}

    \item \textbf{Personal Reflection (Due Saturday Night)}
            \begin{itemize}
                \item Reflect on your learning process and the feedback received.
                \item Optional: Use a preferred AI model to refine reflections:
                \begin{quote}
                    \texttt{"I've completed a reflection on my learning process for this week's synthesis assignment. Here's my reflection [insert reflection]. Could you provide feedback or ask questions to provoke further thought?"}
                \end{quote}
            \end{itemize}
\end{enumerate}

\subsubsection*{Key Weeks and Capstone Activities}

\paragraph{Week 1: Public Administration Theory I (Asynchronous)}
\begin{itemize}
    \item \textbf{Mode}: Asynchronous (video provided).
    \item \textbf{Content}:
    \begin{itemize}
        \item Video lecture and foundational readings:
        \begin{itemize}
            \item \citet{wilson1887}, *The Study of Administration*.
            \item \citet{Weber1946}, *Bureaucracy*.
            \item \citet{gulick1937}, *Notes on the Theory of Organization*.
            \item \citet{Follett1926}, *The Giving of Orders*.
            \item \citet{simon1946}, *Proverbs of Administration*.
        \end{itemize}
    \end{itemize}
    \item \textbf{Assignments}:
    \begin{itemize}
        \item \textbf{Annotated Bibliography (Optional)}: Students may complete this to prepare for Week 2.
        \item No synthesis or peer review this week.
    \end{itemize}
\end{itemize}

\paragraph{Weeks 2--9: Standard Weekly Workflow}
Each week follows the established structure:
\begin{itemize}
    \item Annotated Bibliography (due Tuesday before class).
    \item In-Person Discussion (Tuesday).
    \item Synthesis Draft (due Thursday).
    \item Peer Review (due Friday).
    \item Final Synthesis \& Reflection (due Saturday).
\end{itemize}

\paragraph{Week 10: Capstone Synthesis}
\begin{itemize}
    \item \textbf{In-Person Class (Tuesday)}: Final synthesis preparation.
    \begin{itemize}
        \item \textbf{Activities}:
        \begin{itemize}
            \item Review key themes from the course.
            \item Peer brainstorming and feedback on capstone ideas.
        \end{itemize}
    \end{itemize}
    \item \textbf{Capstone Synthesis Paper}:
    \begin{itemize}
        \item \textbf{Objective}: Integrate course material into a comprehensive analysis.
        \item \textbf{Requirements}:
        \begin{itemize}
            \item Length: 8–10 pages.
            \item Synthesize weekly topics into a cohesive argument.
            \item Highlight themes, patterns, contradictions, and future implications.
        \end{itemize}
        \item \textbf{Assignment}: Final Capstone Submission (due Saturday night)
    \end{itemize}
\end{itemize}

\subsubsection*{Grading Breakdown}
\begin{table}[h!]
\centering
\begin{tabular}{|l|c|}
\hline
\textbf{Component} & \textbf{Weight (\%)} \\
\hline
Weekly Annotated Bibliographies & 20\% \\
Weekly Synthesis Papers & 40\% \\
Peer Reviews & 15\% \\
Capstone Synthesis Paper & 20\% \\
Reflections & 5\% \\
\hline
\end{tabular}
\end{table}

\paragraph{} Note: Students may skip one week of the annotated bibliography, synthesis, and reflection assignment during the semester. This option is intended to accommodate unexpected events or personal commitments. If you choose to skip a week, please notify the professor in advance. In place of the final synthesis during the skip week place a note in the Canvas assignment indicating that you are using your skip week.

\subsection*{2) Reading Discussion Facilitation}
Students will be responsible for facilitating the discussion of the assigned readings each week, for Weeks 3 through 9. This can be done individually or in pairs. The facilitators will lead the class in a discussion of the readings, pose questions to stimulate critical thinking and analysis, and encourage active participation from all students. The facilitators will also provide a brief summary of the readings and highlight the key points for discussion. The facilitators will be selected at the beginning of the semester, and the schedule will be posted on \emph{Canvas}. The facilitators will be graded based on the following criteria:

\begin{itemize}
    \item Preparation: The facilitators are well-prepared and have a thorough understanding of the readings.
    \item Engagement: The facilitators actively engage the class in discussion and encourage participation from all students.
    \item Critical Thinking: The facilitators pose thought-provoking questions and stimulate critical thinking and analysis.
    \item Communication: The facilitators communicate clearly and effectively, ensuring that all students understand the key points of the readings.
    \item Leadership: The facilitators lead the discussion in a professional and respectful manner, creating a positive learning environment.
    \item Participation: The facilitators actively participate in the discussion and contribute to the class's understanding of the readings.
\end{itemize}

\subsection*{3) MPA Comprehensive General Area Essay Exam}

Students will complete a comprehensive general area essay exam as part of the MPA program's comprehensive exam requirement. The exam will consist of two questions from which students will choose one to answer. The questions will be based on the course readings and discussions and will require students to demonstrate their understanding of public administration's key concepts, theories, and debates. The exam will allow students to synthesize their learning in the course and demonstrate their ability to think critically and write clearly about complex issues in public administration. Students who do not pass the exam on the first attempt will have the opportunity to retake the exam once during finals week. The grade for the exam is pass/fail. Students who do not pass on the second attempt will be required to retake the course.

\subsection*{4) Concentration Area Writing Project}

\subsubsection*{Assignment Overview}

This capstone project allows students to deepen their understanding of their MPA concentration area by integrating theoretical frameworks with real-world challenges. Using \textbf{Dr. Mona Hanna-Attisha's \textit{What the Eyes Don’t See} (2018)} as a central text, students will explore how the Flint water crisis exemplifies issues in their concentration, such as public finance, human resources, local government management, or public policy. The assignment is scaffolded into stages to guide students toward producing a polished, analytical final paper.

\paragraph*{Objectives:}
\begin{itemize}
    \item Develop expertise in a specific public administration concentration area.
    \item Enhance research, analytical, and writing skills.
    \item Synthesize theoretical frameworks with practical case studies.
    \item Critically examine the Flint water crisis and its implications for public administration.
\end{itemize}

\paragraph*{Assignment Stages and Requirements:}
\begin{enumerate}
    \item \textbf{Topic Selection (10\% of assignment grade):}
    \begin{itemize}
        \item Submit a one-page proposal outlining your topic, research question, and how you will integrate \textit{What the Eyes Don’t See} into your analysis.
    \end{itemize}
    \item \textbf{Literature Review (40\% of assignment grade):}
    \begin{itemize}
        \item Draft a literature review that synthesizes key themes from \textit{What the Eyes Don’t See}, course readings, and additional scholarly materials.
        \item Discuss how the Flint water crisis ties into themes relevant to your concentration area.
    \end{itemize}
    \item \textbf{Final Paper (50\% of assignment grade):}
    \begin{itemize}
        \item Submit a polished final paper that integrates feedback from earlier stages.
        \item Analyze the Flint water crisis through the lens of your concentration, connecting theoretical frameworks to practical lessons.
        \item Ensure the paper is 12-15 pages (excluding references), double-spaced, and formatted in APA or Chicago author-date style.
    \end{itemize}
\end{enumerate}

\paragraph*{Submission Timeline:}
\begin{itemize}
    \item \textbf{Topic Outline:} Due Week 13
    \item \textbf{Literature Review:} Due Week 14
    \item \textbf{Final Paper:} Due Week 16
\end{itemize}

\paragraph*{Grading Criteria:}
\begin{itemize}
    \item \textbf{Depth of Analysis (30\%):} Engagement with \textit{What the Eyes Don’t See} and the broader literature.
    \item \textbf{Use of Sources (20\%):} Relevance and critical evaluation of sources.
    \item \textbf{Clarity and Organization (20\%):} Logical structure and coherent arguments.
    \item \textbf{Writing Quality (15\%):} Grammar, style, and proper formatting.
    \item \textbf{Completeness and Accuracy (15\%):} Adherence to assignment requirements and deadlines.
\end{itemize}


\section{Course Requirements Due Dates}

The due dates for the course requirements are as follows:
    \begin{itemize}
        \item Annotated Bibliography and Synthesis: 
        \begin{itemize}
            \item Annotated Bibliographies and Syntheses: Due each week by 7:00 p.m. on Tuesday
            \item Peer Review: Due each week by 11:59 p.m. on Friday
            \item Final Synthesis: Due each week by 11:59 p.m. on Saturday
            \item Personal Reflection: Due each week by 11:59 p.m. on Saturday
        \end{itemize}
        \item MPA Comprehensive General Area Essay Exam:
        \begin{itemize}
            \item Distributed on Friday, 4/11 at 8:00 a.m.
            \item Due on Thursday, 4/17 by 4:59 p.m.
        \end{itemize}
        \item Concentration Area Paper:
        \begin{itemize}
            \item Topic Selection: Due on 4/26
            \item Literature Review Draft: Due on 5/3
            \item Final Draft: Due on 5/13 by 9:00 p.m.
        \end{itemize}
    \end{itemize}

\section{Grades}


\subsection*{Grading Scale and Grade Weights}  

The grading scale is shown in Table~\ref{tab:grading-scale}. Grades will be given based on the weights in Table~\ref{tab:grade-weights}.

\begin{table}[h]
\centering
\caption{Grading Scale}
\begin{tabular}{llll}
\toprule
\textbf{Grade} & \textbf{Percentage} & \textbf{Grade} & \textbf{Percentage} \\
\midrule
A+ & 98.0 -- 100 & B- & 80.0 -- 81.9\\
A & 92.0 -- 97.9 & C+ & 78.0 -- 79.9\\
A- & 90.0 -- 91.9 & C & 72.0 -- 77.9\\
B+ & 88.0 -- 89.9 & C- & 70.0 -- 71.9\\
B & 82.0 -- 87.9 & D & 60.0 -- 69.9\\
D- & 50.0 -- 59.9 & F & 0 -- 49.9\\

\bottomrule
\end{tabular}
\label{tab:grading-scale}
\end{table}

\begin{table}[h!]
    \centering
    \caption{Grade Weights}
    \begin{tabular}{ll}
        \toprule
    \textbf{Assignment} & \textbf{Percentage} \\
    \midrule
    Weekly Readings Assignments & 60\% \\
    Reading Discussion Facilitation & 5\% \\
    MPA Comprehensive General Area Essay Exam & 35\% \\
    Concentration Area Paper & 30\% \\
    \bottomrule
    \end{tabular}
    \label{tab:grade-weights}
    \end{table}
    
\section{Course Schedule}

\subsection*{Week 1, Starting 1/19: Introduction and Syllabus Review}
\begin{itemize}
    \item \textbf{Asynchronous Online Session}: Introduction to the Course and Syllabus Review
    \item Readings:
        \begin{itemize}
            \item Review the course syllabus
        \end{itemize}
    \item Assignments:
        \begin{itemize}
            \item Read and annotate the syllabus
            \item Prepare for Week 2 readings
        \end{itemize}
\end{itemize}
\subsection*{Week 2, Starting 1/26: Public Administration Theory I}
\begin{itemize}
    \item \textbf{In-person Session}: Introduction to the Course
    \item Readings:
        \begin{itemize}
            \item \citet{wilson1887}, ``The Study of Administration''
            \item \citet{Weber1946}, ``Bureaucracy'' 
            \item \citet{gulick1937}, ``Notes on the Theory of Organization'' 
            \item \citet{Follett1926}, ``The Giving of Orders'' 
            \item \citet{simon1946}, ``Proverbs of Administration'' 
        \end{itemize}
    \item Due Tuesday: Annotated Bibliography
    \item Due Thursday: Synthesis Paper
    \item Due Friday: Peer Review
    \item Due Saturday: Final Synthesis and Personal Reflection
\end{itemize}

\subsection*{Week 3, Starting 2/2: Ethics and Values in Public Administration}
\begin{itemize}
    \item \textbf{In-person Session}: Public Service Values and Ethics
    \item Readings:
        \begin{itemize}
            \item \citet{friedrich1935}, ``Responsible Government Service Under the American Constitution'' 
            \item \citet{FINER1941}, ``Administrative Responsibility in Democratic Government'' 
            \item \citet{goss1996}, ``A Distinct Public Administration Ethics?'' 
            \item \citet{Adams2009}, ``Unmasking Administrative Evil'' 
            \item \citet{Denhardt2015}, \emph{The New Public Service}, Chapter 7 
        \end{itemize}
    \item Due Tuesday: Annotated Bibliography
    \item Due Thursday: Synthesis Paper
    \item Due Friday: Peer Review
    \item Due Saturday: Final Synthesis and Personal Reflection
\end{itemize}

\subsection*{Week 4, Starting 2/9: Public Administration Theory II}
\begin{itemize}
    \item \textbf{In-person Session}: Public Administration in the U.S. Context
    \item Readings:
        \begin{itemize}
            \item \citet{Stivers2000}, \emph{Bureau Men, Settlement Women} 
            \item \citet{Denhardt2015}, \emph{The New Public Service}, Chapters 1--4 
        \end{itemize}
    \item Due Tuesday: Annotated Bibliography
    \item Due Thursday: Synthesis Paper
    \item Due Friday: Peer Review
    \item Due Saturday: Final Synthesis and Personal Reflection
\end{itemize}

\subsection*{Week 5, Starting 2/16: Leadership and Motivation}
\begin{itemize}
    \item \textbf{In-person Session}: Leadership and Motivation
    \item Readings:
        \begin{itemize}
            \item \citet{Christensen2017}, ``Public Service Motivation Research'' 
            \item \citet{Denhardt2015}, \emph{The New Public Service}, Chapter 8 
            \item \citet{Lachance2017}, ``Public Service Motivation'' 
            \item \citet{maslow1943}, ``A Theory of Human Motivation'' 
            \item \citet{Fairholm2004}, ``Different Perspectives on the Practice of Leadership'' 
        \end{itemize}
    \item Due Tuesday: Annotated Bibliography
    \item Due Thursday: Synthesis Paper
    \item Due Friday: Peer Review
    \item Due Saturday: Final Synthesis and Personal Reflection
\end{itemize}

\subsection*{Week 6, Starting 2/23: Performance Management}
\begin{itemize}
    \item \textbf{In-person Session}: Performance Management
    \item Readings:
        \begin{itemize}
            \item \citet{Behn2003}, ``Why Measure Performance?'' 
            \item \citet{Denhardt2015}, \emph{The New Public Service}, Chapter 9 
            \item \citet{douglas2021}, ``Getting a Grip on Performance of Collaborations'' 
            \item \citet{marvel2015}, ``Unconscious Bias in Citizens' Evaluations\dots'' 
            \item \citet{nicholson-crotty2004}, ``Public Management and Organizational Performance'' 
        \end{itemize}
    \item Due Tuesday: Annotated Bibliography
    \item Due Thursday: Synthesis Paper
    \item Due Friday: Peer Review
    \item Due Saturday: Final Synthesis and Personal Reflection
\end{itemize}

\subsection*{Week 7, Starting 3/2: Street-Level Bureaucrats}
\begin{itemize}
    \item \textbf{In-person Session}: Street-Level Bureaucrats
    \item Readings:
        \begin{itemize}
            \item \citet{Lipsky2010}, \emph{Street-Level Bureaucracy}, Chapters 1--10 
            \item \citet{Denhardt2015}, \emph{The New Public Service}, Chapters 5--6 
        \end{itemize}
    \item Due Tuesday: Annotated Bibliography
    \item Due Thursday: Synthesis Paper
    \item Due Friday: Peer Review
    \item Due Saturday: Final Synthesis and Personal Reflection
\end{itemize}

\subsection*{Week 8, Starting 3/9: Privatization and Contracting}
\begin{itemize}
    \item \textbf{In-person Session}: Privatization and Contracting
    \item Readings:
        \begin{itemize}
            \item \citet{MILWARD2000a}, ``Governing the Hollow State''
            \item \cite{hood1991}, ``A Public Management for All Seasons?''
            \item \citet{brownManagingPublicService2006}, ``Managing Public Service Contracts''
            \item \citet{jos2009}, ``Keeping it Public''
            \item \citet{raineyComparingPublicPrivate2000}, ``Comparing Public and Private Organizations''
        \end{itemize}
    \item Due Tuesday: Annotated Bibliography
    \item Due Thursday: Synthesis Paper
    \item Due Friday: Peer Review
    \item Due Saturday: Final Synthesis and Personal Reflection
\end{itemize}

\subsection*{Week 9, Starting 3/16: 21st Century Challenges}
\begin{itemize}
    \item \textbf{In-person Session}: 21st Century Challenges
    \item Readings:
        \begin{itemize}
            \item \citet{maynard-moody2012}, ``Social Equities and Inequities in Practice'' 
            \item \citet{GOODEN2017}, ``Social Equity and Evidence'' 
            \item \citet{mccandless2022}, ``A Long Road''
            \item \citet{Denhardt2015}, \emph{The New Public Service}, Chapters 10--12 
        \end{itemize}
    \item Due Tuesday: Annotated Bibliography
    \item Due Thursday: Synthesis Paper
    \item Due Friday: Peer Review
    \item Due Saturday: Final Synthesis and Personal Reflection
\end{itemize}

\subsection*{Week 10, Starting 3/23: Capstone Synthesis Paper}
\begin{itemize}
    \item \textbf{In-person Session}: Capstone Synthesis
    \item Readings:
        \begin{itemize}
            \item Review key themes from the course
            \item Peer brainstorming and feedback on capstone ideas
        \end{itemize}
    \item Due Saturday: Final Capstone Synthesis Paper Submission
\end{itemize}

\subsection*{Spring Break, Starting 3/30}
\begin{itemize}
    \item \textbf{No Classes}: Enjoy your break!
\end{itemize}

\subsection*{Week 11, Starting 4/6: Comprehensive General Area Essay Exam}
\begin{itemize}
    \item \textbf{Asynchronous Session}: Comprehensive General Area Essay Exam Distributed on Friday at 8:00 a.m.
\end{itemize}

\subsection*{Week 12, Starting 4/13: Comprehensive General Area Essay Exam}
\begin{itemize}
    \item \textbf{Asynchronous Session}: Comprehensive General Area Essay Exam Due on Thursday at 4:59 p.m.
\end{itemize}

\subsection*{Week 13, Starting 4/20: Concentration Area Paper}
\begin{itemize}
    \item \textbf{Asynchronous Session}: Concentration Area Paper
    \item Paper expectations and guidelines 
    \item Schedule a 1:1 meeting to discuss topic selection at \href{https://dadams.io/appointments}{dadams.io/appointments}
    \item Host 1:1 meetings on 4/22 at 7:00 - 9:00 p.m. at \href{https://fullerton.zoom.us/j/3347502639}{fullerton.zoom.us/j/3347502639}
    \item Due: Topic Selection
\end{itemize}

\subsection*{Week 14, Starting 4/27: Concentration Area Paper}
\begin{itemize}
    \item \textbf{Asynchronous Session}: Concentration Area Paper
    \item Due: Literature Review Draft
\end{itemize}

\subsection*{Week 15, Starting 5/4: Concentration Area Paper}
\begin{itemize}
    \item \textbf{In-Person Session}
    \item Course Wrap-Up and Final Reflections
    \item Potluck Celebration
\end{itemize}

\subsection*{Week 16, Starting 5/11: Concentration Area Paper}
\begin{itemize}
    \item \textbf{Asynchronous Session}: Concentration Area Paper
    \item \textbf{Due: Concentration Area Paper Final Draft on 5/13 by 9:OO p.m.}
\end{itemize}

            \singlespace
            \bibliographystyle{apsr}
            \bibliography{521_students}

% updated date and time
\begin{flushright}
Updated: \today
\end{flushright}

            
\end{document}